% Transfer of Status qualifying dissertation.
% Structure: Introduction; literature review (lit_review.tex);
% Project Proposal (proposal.tex). Shared preamble in preamble.tex.
% =====================================================================
% Shared preamble for draft.tex, draft-communication.tex, draft-interventional.tex
% Orange [FILL: ...] marks facts only you can supply. See skeleton.tex
% section 0.5 for the qualifiers that must survive editing.
% =====================================================================
\documentclass[11pt,a4paper]{article}
\usepackage[margin=2.2cm]{geometry}
\usepackage[backend=biber,style=numeric,sorting=nyt,maxbibnames=6,giveninits=true]{biblatex}
\usepackage{amsmath,amssymb,amsthm}
% \usepackage{amsthm}
% \usepackage[margin=2cm]{geometry}
\usepackage{dsfont}
\usepackage{tikz}
\usetikzlibrary{cd}
\usetikzlibrary{babel,calc}
\usetikzlibrary{positioning,arrows.meta,fit,backgrounds}
% \usepackage{tikz-cd}
% \usepackage{titlesec}
\usepackage{enumitem}
\usepackage{hyperref}
\usepackage{cleveref}
% \usepackage{parskip}
\usepackage{ebproof}
% \usepackage{quiver}
\usepackage{stackengine,scalerel}
% \usepackage{xparse}
% \usepackage{physics}
% \usepackage{physics2}
% \usephysicsmodule{ab,ab.braket}
% \usepackage{tikzit}
% \input{smc.tikzstyles}
\usepackage{xcolor}
\usepackage{graphicx}
\usepackage{array}    % m{} columns + \arraybackslash (combined car-model table)
\usepackage{booktabs} % \toprule/\midrule/\bottomrule

\usepackage{mathtools}
\usepackage{stmaryrd}
\usepackage{subcaption}

\newtheorem*{fact}{Fact}

\crefname{prop}{prop.}{props.}
\Crefname{prop}{Proposition}{Propositions} % Capitalized version for start of sentences

\crefname{remark}{remark}{remarks}
\Crefname{remark}{Remark}{Remarks} % Capitalized version for start of sentences

\crefname{appendix}{appendix}{appendices}
\Crefname{appendix}{Appendix}{Appendices} % Capitalized version for start of sentences

% \crefname{property}{property}{properties}
% \Crefname{property}{Property}{Properties} % Capitalized version for start of sentences

\crefformat{property}{#2#1#3}
\Crefformat{property}{#2#1#3}

\crefname{def}{definition}{definitions.}
\Crefname{def}{Definition}{Definitions} % Capitalized version for start of sentences

\crefformat{section}{#2\S#1#3}
\Crefformat{section}{#2\S#1#3}

\usepackage{listings}

\usepackage{tcolorbox}
\tcbuselibrary{listings} % Import listings support

\definecolor{codegreen}{rgb}{0,0.6,0}
\definecolor{codegray}{rgb}{0.5,0.5,0.5}
\definecolor{codepurple}{rgb}{0.58,0,0.82}
\definecolor{backcolour}{rgb}{0.95,0.95,0.92}

\lstdefinestyle{mystyle}{
    backgroundcolor=\color{backcolour},
    commentstyle=\color{codegreen}\textit,
    keywordstyle=\color{blue},
    numberstyle=\tiny\color{codegray},
    stringstyle=\color{codepurple},
    basicstyle=\ttfamily\footnotesize,
    breakatwhitespace=false,
    breaklines=false,
    captionpos=b,
    keepspaces=true,
    numbers=left,
    numbersep=5pt,
    showspaces=false,
    showstringspaces=false,
    showtabs=false,
    tabsize=2
}

\lstset{style=mystyle}

\tcbset{
    mystyle/.style={
        listing only,
        breakable, % Allows box to span pages
        colback=CodeBack,
        colframe=CodeFrame,
        boxrule=0.5pt,
        sharp corners,
        listing options={
            backgroundcolor=\color{backcolour},   
            commentstyle=\color{codegreen}\textit,
            keywordstyle=\color{blue},
            numberstyle=\tiny\color{codegray},
            stringstyle=\color{codepurple},
            basicstyle=\ttfamily\footnotesize,
            breakatwhitespace=false,         
            breaklines=true,                 
            captionpos=b,                    
            keepspaces=true,                 
            numbers=left,                    
            numbersep=5pt,                  
            showspaces=false,                
            showstringspaces=false,
            showtabs=false,                  
            tabsize=2,
        }
    }
}

\newtcblisting{codelisting}[3][]{
    listing only,
    % breakable, % Allows box to span pages
    colback=backcolour,
    % colframe=CodeFrame,
    boxrule=0pt,
    boxsep=0pt,
    top=1mm,
    bottom=1mm,
    sharp corners,
    notitle,
    listing options={
        language=#2,
        morekeywords= {exo, endo},
        commentstyle=\color{codegreen}\textit,
        keywordstyle=\color{blue},
        numberstyle=\tiny\color{codegray},
        stringstyle=\color{codepurple},
        basicstyle=\ttfamily\tiny,
        breakatwhitespace=false,         
        breaklines=true,                 
        captionpos=b,                    
        keepspaces=true,                 
        numbers=left,                    
        numbersep=5pt,                  
        showspaces=false,                
        showstringspaces=false,
        showtabs=false,                  
        tabsize=2,
        aboveskip=0pt,  % Removes the space above the code
        belowskip=0pt,  % Removes the space below the code
        mathescape=#3,
        #1               % optional extra listings options (e.g. escapeinside for \cnew)
    }
}

% Highlight code that is new relative to the previous rung's program.
% Use inside a codelisting that passes escapeinside, e.g.
%   \begin{codelisting}[escapeinside={(*}{*)}]{haskell}{true} ... (*\cnew{new}*) ...
\definecolor{codenew}{RGB}{255,228,130}
\newcommand{\cnew}[1]{{\setlength{\fboxsep}{0.6pt}\colorbox{codenew}{\ttfamily #1}}}
% Coloured (not highlighted) font, for marking new syntax in the grammar figure.
\newcommand{\cnewm}[1]{\textcolor{blue!65!black}{#1}}

% Theo's Macros
\newcommand{\naturals}{\mathbb{N}}
\newcommand{\integers}{\mathbb{Z}}
\newcommand{\booleans}{\mathbb{B}}
\newcommand{\indicator}{\mathds{1}}
\newcommand{\va}{\mathbf{a}}
\newcommand{\vb}{\mathbf{b}}
\newcommand{\vx}{\mathbf{x}}
\newcommand{\vy}{\mathbf{y}}
\newcommand{\vw}{\mathbf{w}}
\newcommand{\vz}{\mathbf{z}}
\newcommand{\vsigma}{\mathbf{\sigma}}
\newcommand{\reals}{\mathbb{R}}
\newcommand{\complex}{\mathbb{C}}
\newcommand{\field}{\mathbb{F}}
\newcommand{\sign}{\mathrm{sign}}
\DeclareMathOperator*{\argmin}{arg\,min}
\newcommand{\true}{\mathsf{True}}
\newcommand{\false}{\mathsf{False}}
\newcommand{\brackets}[1]{\{#1\}}

% \newcommand{\bicatb}{\mathcal{B}}

\newcommand{\catc}{\mathbf{C}}
\newcommand{\catcop}{\catc^{op}}
\newcommand{\catd}{\mathbf{D}}
\newcommand{\cate}{\mathcal{E}}
\newcommand{\eqdef}{\triangleq}
\newcommand{\id}{\mathbf{id}}
% \newcommand{\idbicat}{\mathbf{Id}}
\newcommand{\Id}{\mathbf{Id}}
\newcommand{\righttop}[1]{\xrightarrow{#1}}
\newcommand{\lefttop}[1]{\xleftarrow{#1}}
\newcommand{\Ob}{\mathbf{Ob} \ }
\newcommand{\Hom}{\mathrm{Hom}}
% \newcommand{\yo}{\textmc{よ}}
\newcommand{\yo}{\textbf{y}}

\newcommand{\terminalmorph}{\langle\rangle}
\newcommand{\prodmorph}[2]{\langle#1, #2\rangle}

\newcommand{\ul}[1]{\underline{#1}}
\newcommand{\ulyo}{\ul{\yo}}
\newcommand{\yozero}{\yo(0)}
\newcommand{\yoone}{\yo(1)}
\newcommand{\yon}{\yo(n)}
\newcommand{\yom}{\yo({m})}
\newcommand{\yop}{\yo(p)}
% \newcommand{\op}{\mathbf{op}}
\newcommand{\Ops}{\mathcal{O}}
% \newcommand{\typing}[3]{#1 \ | \ #2 \vdash #3}
% \newcommand{\typingfilled}[1]{\Gamma \ | \ \Delta \vdash #1}
% \newcommand\eye{\ensurestackMath{\stackinset{c}{}{c}{-.33pt}%
%   {\bullet}{\bigcirc}}}
% \newcommand\smallereye{\scalerel*{\eye}{x}}
% \newcommand{\eyemid}{ \ \eye \ }

% Vectors
\newcommand{\veca}{\Vec{a}}
\newcommand{\vecb}{\Vec{b}}
\newcommand{\vecc}{\Vec{c}}
\newcommand{\vecd}{\Vec{d}}
\newcommand{\vecx}{\Vec{x}}
\newcommand{\vecy}{\Vec{y}}
\newcommand{\vecz}{\Vec{z}}

\newcommand{\alphaequiv}{\simeq_\alpha}

\newcommand{\LHS}{\mathsf{LHS}}
\newcommand{\RHS}{\mathsf{RHS}}

% Chapter 2 algebraic theories
\newcommand{\Mon}{\mathsf{Mon}}

\newcommand{\Tree}{\mathsf{Tree}}

\newcommand{\Kl}[1]{\mathsf{Kl}({#1})}

\newcommand{\Name}{\mathcal{N}}
\newcommand{\Var}{\mathsf{Var}}
\newcommand{\IPVar}{\mathsf{IPVar}}
\newcommand{\GroundValue}{\mathsf{GroundValue}}
\newcommand{\Value}{\mathsf{Value}}
\newcommand{\Expr}{\mathsf{Expr}}

\newcommand{\bind}{\mathbin{>\!\!\!>\mkern-6.7mu=}}
\newcommand{\fin}{\mathsf{fin}}
\newcommand{\MVal}{\mathsf{MVal}}
\newcommand{\Intervention}{\mathsf{Intv}}
\newcommand{\MultiWorld}{\mathsf{MultiWorldCounterfactual}}
\newcommand{\TwinWorld}{\mathsf{TwinWorldCounterfactual}}
\newcommand{\supp}{\mathsf{supp}}

% ------------- Formal Def ----------------
\newcommand{\IdealisedChiRho}{\mathsf{ChiRho}}
\newcommand{\Prob}{\mathsf{Prob}}
\newcommand{\Dist}{\mathsf{Dist}}
\newcommand{\Paths}{\mathsf{Paths}}
\newcommand{\UPaths}{\mathsf{UPaths}}
\newcommand{\posreals}{\reals_{\geq 0}}

% Syntax
% \newcommand{\opif}[3]{\mathbf{if} \ #1 \ #2 \ #3}
% \newcommand{\lift}[1]{\overline{#1}}
\newcommand{\opprod}{\mathsf{prod}}
\newcommand{\opscore}{\mathsf{score}}
\newcommand{\opflip}{\mathsf{flip}}
\newcommand{\dbernoulli}{\mathbf{bernoulli}}
\newcommand{\dnormal}{\mathbf{normal}}
\newcommand{\constr}{\mathbf{constr}}
\newcommand{\intvnone}{\mathsf{None}}
\newcommand{\intvsome}{\mathsf{Some}}
\newcommand{\opskip}{\mathsf{skip}}
\newcommand{\opsample}{\mathsf{sample}}
\newcommand{\opintervene}{\mathsf{intervene}}
\newcommand{\opbroadcast}{\mathsf{broadcast}}
\newcommand{\optrack}{\mathsf{track}}
\newcommand{\opdistribute}{\mathsf{distribute}}
\newcommand{\opif}{\mathbf{if}}
\newcommand{\ifelse}[3]{\opif \ #1 \ \mathbf{then} \ #2 \ \mathbf{else} \ #3}
\newcommand{\opthunk}{\mathsf{thunk}}
\newcommand{\opforce}{\mathsf{force}}
\newcommand{\returnSymb}{\mathbf{return}}
\newcommand{\return}[1]{\returnSymb \ #1}
\newcommand{\returnParam}[2]{\returnSymb_{#1} \ #2}
\newcommand{\app}{\odot}
\newcommand{\pureval}{\mathbf{pureVal}}
\newcommand{\mem}{\mathsf{mem}}
\newcommand{\letbind}[3]{\mathbf{let} \ #1 = #2 \ \mathbf{in} \ #3}
\newcommand{\opdo}{\mathsf{do}}
\newcommand{\rhide}{\mathsf{hide}}
\newcommand{\rwk}{\mathsf{wk}}
\newcommand{\rcontr}{\mathsf{contr}}
\newcommand{\rexch}{\mathsf{exch}}
\newcommand{\rrename}{\mathsf{rename}}
\newcommand{\Renamings}{\mathsf{Renamings}}
\newcommand{\Deltahat}{\widehat{\Delta}}
\newcommand{\projf}{\pi_{\mathsf{f}}}
\newcommand{\projcf}{\pi_{\mathsf{cf}}}
\newcommand{\coerce}[3]{\mathsf{coerce}_{#2}^{#1}[#3]}
\newcommand{\rv}{\underline{r}}
\newcommand{\restr}{\rvert}
\newcommand{\subst}[2]{[#1/#2]}
\newcommand{\ket}[1]{| #1 \rangle}

% Types
\newcommand{\tnat}{\mathsf{Nat}}
\newcommand{\tbool}{\mathsf{Bool}}
\newcommand{\tint}{\mathsf{Int}}
\newcommand{\tposreal}{\mathsf{PosReal}}
\newcommand{\treal}{\mathsf{Real}}
\newcommand{\tunit}{\mathsf{Unit}}
\newcommand{\tprod}{\mathsf{Prod}}
\newcommand{\tP}[2]{P \ #1 \ #2}
\newcommand{\unitgrade}{[]}
\newcommand{\rarrow}{\Rightarrow}

% Type system
\newcommand{\dtyping}[3]{#1 \vdash^d #2: #3}
\newcommand{\ptyping}[4]{#1 \mid #2 \vdash^p #3: #4}
\newcommand{\rtyping}[4]{#1 \vdash #2: #3 \rarrow #4}

\newcommand{\ctxequiv}{\simeq_{\mathsf{ctx}}}
\newcommand{\Values}{\mathsf{Values}}

% DenSem
% \newcommand{\sem}[1]{\llbracket #1 \rrbracket}
\newcommand{\subfin}{\subseteq_{\fin}}
\newcommand{\inj}{\mathbf{inj}}
\newcommand{\detsem}[1]{[#1]}
\newcommand{\List}{\mathrm{List}}
\newcommand{\sIntervention}{\mathrm{Intv}}
\newcommand{\sP}[2]{\mathbf{P} \ #1 \ #2}
\newcommand{\sPP}{\mathbf{P}}
\newcommand{\sVal}{\mathbf{Val}}
\newcommand{\gradecatsymb}{\mathbb{G}}
\newcommand{\gradecat}[1]{\gradecatsymb_{#1}}
\newcommand{\sdist}{\mathrm{dist}}
\newcommand{\sintervene}{\mathrm{intervene}}

\DeclareMathOperator*{\dom}{dom}

% New language
\newcommand{\sCaus}{\mathrm{Caus}}
\newcommand{\Caus}[2]{\sCaus \ #1 \ #2}
% \newcommand{\Maybe}{\mathsf{Maybe}}
\newcommand{\opnew}{\mathsf{new}}
\newcommand{\Caustic}{\mathsf{\lambda_C}}
\newcommand{\cdtyping}[3]{#1 \vdash^d_{\Caustic} #2: #3}
\newcommand{\crtyping}[4]{#1 \vdash^r_{\Caustic} #2: #3 \rarrow #4}
\newcommand{\cptyping}[4]{#1 \vdash^{p}_{\Caustic} #2 :#3 \& #4}
\newcommand{\cptypingreordered}[4]{\cptyping{#2}{#3}{#4}{#1}}

% Compilation and Logical Relation
\DeclareMathOperator*{\len}{dom}

\newcommand{\Maybe}[1]{\mathrm{Maybe}({#1})}
\newcommand{\sem}[1]{\llbracket{#1}\rrbracket}
\newcommand{\Just}[1]{\mathrm{Just}({#1})}
\newcommand{\Nothing}{\mathrm{Nothing}}

% Dario's Macros
\usepackage{bm}
\renewcommand{\vec}{\bm}

\newcommand{\todo}[1]{\textcolor{blue}{[TODO: #1]}}
\newcommand{\prompt}[1]{\textcolor{purple}{[PROMPT: #1]}}

\newcommand{\cat}[1]{\mathbb{#1}}
\newcommand{\catname}[1]{\mathbf{#1}}

\newcommand{\set}{\mathbf{set}}
\newcommand{\Set}{\catname{Set}}

\newcommand{\concat}{+\!\!+}

\newcommand{\defeq}{\coloneq}


\newcommand{\C}{\cat C}
\newcommand{\G}{\mathcal G}

\newcommand{\op}{\mathrm{op}}


\newcommand{\thought}[3]{\marginpar{\tiny\itshape\raggedright\textcolor{#2}{\textbf{#1:} #3}}\ignorespaces} \marginparwidth=1.4cm \marginparsep=0cm
\newcommand{\jeremy}[1]{\thought{Jeremy}{green!20!black}{#1}}

\newcommand{\IntvInstr}{\mathsf{IntvInstr}}
\newcommand{\oplookup}{\mathsf{lookup}}
\newcommand{\oprestrict}{\mathsf{restrict}}
\newcommand{\FinSet}{\catname{FinSet}}
\newcommand{\Bij}{\catname{Bij}}
\newcommand{\Type}{\mathbf{Type}}

\newcommand{\vpure}{\mathbf{pure}}
\newcommand{\vapp}{\mathbin{\langle * \rangle}}
% \DeclareMathOperator*{\vapp}{\langle * \rangle}
\newcommand{\fmap}{\mathbin{\langle \$ \rangle}}

\newcommand{\colorone}[1]{\textcolor{purple}{#1}}
\newcommand{\colortwo}[1]{\textcolor{blue}{#1}}
\newcommand{\ParIndex}{\mathsf{ParIndex}}
\newcommand{\Env}{\mathbf{Env}}

\newcommand{\inl}{\mathbf{inl}}
\newcommand{\inr}{\mathbf{inr}}

\newcommand{\note}{\textcolor{purple}}
\newcommand{\cur}{\mathbf{cur}}
\newcommand{\uncur}{\mathbf{uncur}}

\newcommand{\intvPt}{\mathsf{intvPt}}
\newcommand{\Endo}{\mathcal{E}}
\newcommand{\sEndo}{\mathsf{Endo}}

\newcommand{\arrow}{\rightsquigarrow}
\newcommand{\arrarr}{\mathsf{arr}}
\newcommand{\arrcomp}{\ggg}
\newcommand{\arrfirst}{\mathsf{first}}


\newcommand{\DarioMonad}{\mathsf{T}_{\mathsf{CBN}}}
\newcommand{\oprun}{\mathsf{run}}
\newcommand{\dolet}{\mathbf{let}}
\newcommand{\doletexpr}[3]{\dolet \ #1 \leftarrow #2 \ \doin \ #3}
\newcommand{\doletexprParam}[4]{\dolet_{#1} \ #2 \leftarrow #3 \ \doin \ #4}
\newcommand{\doin}{\mathbf{in}}

\newcommand{\typing}[3]{#1 \vdash #2: #3}
\newcommand{\wtyping}[4]{#1 \vdash_{#2} #3: #4}
\newcommand{\polytyping}[3]{#1 \vdash_{[-]} #2: #3}
\newcommand{\TMon}[2]{T_{#1}\left(#2 \right)}

\newcommand{\catIntv}{\mathbf{IntvStruct}}

\newcommand{\bOne}{\mathds{1}}
\newcommand{\newEndo}{\mathsf{endo}}

\newcommand{\Meas}{\mathsf{Meas}}
\newcommand{\perWorld}{\mathsf{perWorld}}
\newcommand{\valueAt}{\mathsf{valueAt}}

\newcommand{\PSh}{\mathbf{PSh}}
\newcommand{\PShBij}{\PSh(\Bij)}

\newcommand{\LangName}{\mathsf{\lambda_C}}
\newcommand{\World}{\mathsf{World}}

\newcommand{\dwtyping}[4]{#1 \mid #2 \vdash^d #3: #4}
\newcommand{\pwtyping}[5]{#1 \mid #2 \mid #3 \vdash^p #4: #5}
% \newcommand{\rtyping}[4]{#1 \vdash #2: #3 \rarrow #4}

\newcommand{\Any}{\mathsf{Any}}
\newcommand{\intvConstr}{\mathsf{intvConstr}}
\newcommand{\twin}{\mathsf{twin}}

\newcommand{\Functions}{\mathsf{Functions}}
\newcommand{\Distributions}{\mathsf{Distributions}}
\newcommand{\Constants}{\mathsf{Constants}}

\newcommand{\rone}[1]{{#1}^{(1)}}
\newcommand{\rtwo}[1]{{#1}^{(2)}}
\newcommand{\rthree}[1]{{#1}^{(3)}}

\newcommand{\Lone}{L_1}
\newcommand{\Ltwo}{L_2}
\newcommand{\Lthree}{L_3}

\newcommand{\Sbs}{\catname{Sbs}}
\newcommand{\Qbs}{\catname{Qbs}}

\newcommand{\typingone}[4]{#1 \vdash_{1}^{#4} #2: #3}
\newcommand{\etypingone}[3]{\typingone{#1}{#2}{#3}{\mathbf{e}}}
\newcommand{\ttypingone}[3]{\typingone{#1}{#2}{#3}{\mathbf{t}}}
\newcommand{\qtypingone}[3]{\typingone{#1}{#2}{#3}{\mathbf{q}}}

\newcommand{\typingtwo}[5]{#1 \mid #2 \vdash_{2}^{#5} #3: #4}
\newcommand{\etypingtwo}[4]{\typingtwo{#1}{#2}{#3}{#4}{\mathbf{e}}}
\newcommand{\ttypingtwo}[4]{\typingtwo{#1}{#2}{#3}{#4}{\mathbf{t}}}
% \newcommand{\qtypingtwo}[4]{\typingtwo{#1}{#2}{#3}{#4}{q}}

\newcommand{\etypingthree}[4]{#1 \mid #2 \vdash_{3}^{\mathbf{e}} #3: #4}
\newcommand{\ttypingthree}[6]{#1; #2 \mid #3; #4 \vdash_{3}^{\mathbf{t}} #5: #6}
% \newcommand{\qtypingtwo}[4]{\typingtwo{#1}{#2}{#3}{#4}{q}}
\newcommand{\Gammaexo}{\Gamma_{\mathsf{exo}}}
\newcommand{\Gammaendo}{\Gamma_{\mathsf{endo}}}
\newcommand{\Deltaexo}{\Delta_{\mathsf{exo}}}
\newcommand{\Deltaendo}{\Delta_{\mathsf{endo}}}

\newcommand{\typingHO}[3]{#1 \vdash_{\mathsf{HO}} #2: #3}


\newcommand{\scoringquery}[4]{\letbind{#1}{#3}{\opscore(#2)};\ \return{#4}}

\newcommand{\xletbind}[3]{\mathbf{let \ exo} \ #1 = #2 \ \mathbf{in} \ #3}
\newcommand{\nletbind}[3]{\mathbf{let \ endo} \ #1 = #2 \ \mathbf{in} \ #3}


\newcommand{\opdefault}{\mathbf{default}}
\newcommand{\opcondition}{\mathsf{condition}}
\newcommand{\opapplyIntv}{\mathsf{applyIntv}}
\newcommand{\forget}{\mathbf{forget}}
\newcommand{\semSCM}[1]{\sem{#1}_{\mathsf{SCM}}}

% Section 4
\newcommand{\IntvConstrName}{\mathbf{IntvC}}
\newcommand{\semone}[1]{\sem{#1}_{\Lone}}
\newcommand{\semtwo}[1]{\sem{#1}_{\Ltwo}}
\newcommand{\semthree}[1]{\sem{#1}_{\Lthree}}
\newcommand{\semho}[1]{\sem{#1}_{\LHO}}
\newcommand{\liftedEndo}{\widehat{\Endo}}
\newcommand{\sliftedEndo}{\sEndo}
\newcommand{\lift}{\mathsf{lift}}
\newcommand{\liftEndo}{\mathsf{endo}}
\newcommand{\catp}{\mathbf{P}}
\newcommand{\App}{\mathsf{App}}
\newcommand{\intvLabel}{\mathsf{intvLabel}}
\newcommand{\ulBij}{\underline{\Bij}}

% Section 5
\newcommand{\LHO}{L_{\mathsf{HO}}}
\newcommand{\LthreeHOcompile}[1]{\llparenthesis #1 \rrparenthesis}
\newcommand{\IntvMonadSymb}{\mathcal{I}}
\newcommand{\IntvMonad}[2]{\IntvMonadSymb_{#1}\left(#2\right)}
\newcommand{\Twinned}{\mathcal{T}}
% \newcommand{\TwinnedIntvT}[1]{\Twinned_{#1} \IntvMonadSymb}
\newcommand{\sTwinned}{\mathsf{Twinned}}
\newcommand{\sIntvMonadSymb}{\mathsf{IntvT}}
\newcommand{\sIntvMonad}[2]{\sIntvMonadSymb_{#1}\left(#2\right)}
\newcommand{\opnewEndo}{\mathsf{newEndo}}
\newcommand{\unzip}{\mathsf{unzip}}
\newcommand{\zip}{\mathsf{zip}}

% commenting
\colorlet{jyColour}{blue}

\makeatletter
\patchcmd{\@@addmarginpar}{\ifodd\c@@page}{\ifodd\c@@page\@@tempcnta\m@@ne}{}{}
\makeatother
\reversemarginpar
\newcommand{\sidenote}[3]{\marginpar{\tiny\itshape\raggedright\textcolor{#2}{\textbf{#1:} #3}}\ignorespaces} \marginparwidth=1.4cm \marginparsep=0cm
\newcommand{\jy}[1]{\sidenote{jy}{jyColour}{#1}}


% SCMs
\DeclarePairedDelimiterX{\parensbar}[1]
{(\!\delimsize\vert}
{\delimsize\vert\!)}
{#1}
\newcommand{\compileSCM}[1]{\ensuremath{\parensbar{#1}}}
\newcommand{\scmIntervene}[2]{\ensuremath \mathsf{do}_{#1, #2}}
\newcommand{\scmTwin}[1]{\ensuremath \mathsf{twin}_{#1}}
\newcommand{\IntClass}{\ensuremath{\mathsf{Intv}}}
\newcommand{\DefaultClass}{\ensuremath{\mathsf{Default}}}

\newcommand{\commentcolor}{\textcolor{codegreen}}
\newcommand{\highlightcolor}{\textcolor{blue}}
\newcommand{\SCM}{\mathcal{M}}

\usepackage{xcolor}
\usepackage{booktabs,array,tabularx}
\usepackage{listings}
\lstset{basicstyle=\ttfamily\small,mathescape=true,columns=fullflexible,keepspaces=true}

\addbibresource{refs.bib}
\addbibresource{refs-quantum.bib}
\addbibresource{refs-new.bib}
\setcounter{biburlnumpenalty}{100}\setcounter{biburlucpenalty}{100}\setcounter{biburllcpenalty}{100}

\newcommand{\fillin}[1]{{\color{orange!80!black}[FILL: #1]}}
\newcommand{\iotq}{\mathsf{IO}\otimes\mathsf{Q}}
\newcommand{\CPinf}{\mathbf{CP}^{\infty}}
\newcommand{\CP}{\mathbf{CP}}
\providecommand{\LHO}{L_{\mathsf{HO}}}
\newcommand{\Lk}[1]{L_{#1}}
\newcommand{\CausC}{\mathbf{Caus}}
\newcommand{\FinStoch}{\mathbf{FinStoch}}
\newcommand{\Mat}{\mathbf{Mat}(\mathbb{R}_{\geq 0})}
\newcommand{\WP}[1]{WP\,#1}
\newcommand{\opin}{\mathsf{in}}
\newcommand{\opout}{\mathsf{out}}
\providecommand{\opnew}{\mathsf{new}}
\newcommand{\opmeasure}{\mathsf{measure}}
\newcommand{\applyU}[1]{\mathsf{apply}_{#1}}
\newcommand{\dok}{\mathsf{do}}
\newtheorem{conj}{Conjecture}
\newtheorem{hyp}{Hypothesis}




\title{Semantics of Causal Probabilistic Programming\\[1ex]
\large Qualifying Dissertation for the Transfer of Status}
\author{Theo Wang\\ \small Department of Computer Science, University of Oxford\\ \small Supervisor: Sam Staton}
\date{Michaelmas Term 2026}

\begin{document}

\maketitle

\begin{abstract}
\noindent
Causal inference asks what would happen if a variable were forced to a value, and the standard answer, Pearl's structural causal models, is a syntactic transformation of a model. Causal probabilistic programming aims to make such interventions primitives of an expressive probabilistic language, and most existing languages do so by evaluation strategies without an account of what they compute. In my first year, with coauthors, I gave that account: a hierarchy of probabilistic languages in which interventions and counterfactual twinning are program transformations, the interventional face of causal structure is defined by a contextual equivalence of programs, and a graded denotational semantics is fully abstract for it over concrete models and extends to higher order. The proposed thesis builds outward from this base in four directions: a practical, vectorised implementation and an extension to cyclic models; a better understanding of the categorical constructions, including a synthetic account in Markov categories; a graph-centred account in which the causal graph rather than the model is primitive, and in which the do-calculus is compared with an equational theory of programs; and, open-endedly, the relation between interventional causality and the notions of causality used in concurrency and in higher-order quantum processes. Sections~\ref{sec:foundations}--\ref{sec:synthesis} review the literature; Section~\ref{sec:proposal} is the proposal.
\end{abstract}

\tableofcontents

\section{Introduction}\label{sec:intro}
% Introduction (input by main.tex after \section{Introduction})

\paragraph{Three notions of causality.}
When a computer scientist says that one part of a system \emph{causes} another, they may mean at least three distinct things. In Pearl's theory of causal inference \cite{pearl2009causality}, $X$ has a causal effect on $Y$ if \emph{intervening} on $X$---forcing it to a value, regardless of how it would otherwise have been produced---changes the distribution of $Y$. In probabilistic modelling, a causal structure is a directed acyclic graph along which a joint distribution \emph{factorises}, and the content of the structure is the family of conditional independences it entails \cite{pearl2009causality,fritz2023dseparation}. In quantum foundations and in the theory of higher-order quantum processes \cite{kissinger2019categorical,oreshkov2012noncausal}, the causal structure of a process is given by its \emph{signalling} constraints, which inputs can influence which outputs, and in concurrency theory a causal structure is an explicit partial order on events \cite{lamport1978time,pratt1986pomsets}.

Three further usages should be named so as to be set aside. In categorical quantum mechanics, a process is called \emph{causal} when it preserves the discarding map \cite{PQP_Coecke_Kissinger_2017}; this is the causality axiom of Chiribella, D'Ariano and Perinotti \cite{chiribella2010purification}, glossed as ``no signalling from the future'', and Coecke \cite{coecke2014terminality} shows that it implies the no-signalling principle for processes composed along a causal structure; it is the base axiom on which signalling structure is built, not a notion of causality between parts. Fritz \cite{fritz2020synthetic} calls a Markov category \emph{causal} when the output of a morphism cannot depend on a choice made downstream of it; this is again a property of the ambient category. In artificial intelligence and in verification, \emph{actual causality} \cite{halpern2016actual} asks which particular event caused which particular outcome in a single run; that question is not treated here. Nor are experimental tests of indefinite causal order or causal discovery from data. Identifiability and the do-calculus are not set aside: the proposal reaches them (\S\ref{sec:prop-graph}).

That the same word carries these meanings is not an accident, but neither is it a theorem. Each notion is \emph{extensional}: it characterises causal structure by what an observer may do to the system---intervene, condition, send signals. What such observers probe is \emph{intensional}: the dataflow of a program, the graph of a model, the wiring of a circuit. Programs carry this intensional structure explicitly in their syntax. The viewpoint of this thesis is therefore that causality is a property of \emph{program semantics}, and that the right way to study a notion of causality is to relate the intensional structure in the syntax to its extensional characterisation in the semantics by a correspondence theorem. A recurring theme, made precise in \S\ref{sec:synthesis}, is that the two never quite coincide: a dependence that is present in the syntax can be invisible to every observer, by what the causal-inference literature calls \emph{fine tuning}.

\paragraph{Where I stand.}
In my first year, with coauthors, I built a hierarchy of first-order probabilistic languages in which Pearl's interventions and counterfactual \emph{twinning} are program transformations, defined program equivalences contextually, and gave a graded denotational semantics that is fully abstract for concrete models and extends to a higher-order language (\S\ref{sec:intv-wherewestand}). This gives the interventional notion its extensional face: equivalence under all interventional contexts, which the graded semantics characterises exactly. What it does not characterise is the dataflow itself, which is strictly finer (\S\ref{sec:intv-wherewestand}). In the running language, a program reads an \emph{intervention label} $\ell$ through $\mathsf{applyIntv}\ \ell\ t$, which yields the label's value when it is set and $t$ otherwise; if a second variable $z$ does not depend on the result, every intervention at $\ell$ leaves the distribution of $z$ unchanged, and the contextual equivalence makes this an equation between programs. Causal independence is \emph{tested by contexts}, and, in programming-language terms, it is a \emph{noninterference} property \cite{goguen1982noninterference}. The thesis takes this work as its subject and builds outward from it: towards practice and expressivity, towards a better understanding of the constructions it rests on, towards a graph-centred account, and, open-endedly, towards the notions of causality used in concurrency and in higher-order quantum processes.

\paragraph{Structure of this document.}
Sections \ref{sec:foundations}--\ref{sec:order} review the literature. Section \ref{sec:foundations} covers the semantic foundations: algebraic effects, probabilistic programs, Markov categories and grading. Section \ref{sec:intv} covers causal models, their equivalences, cyclic models, identifiability, their diagrammatic presentations, and causal probabilistic programming, which is the state of the art the proposal extends. Section \ref{sec:order} positions the two neighbouring fields that the proposal reaches into, higher-order quantum processes and the semantics of concurrency, at the level of their main definitions and results; their detailed survey is a milestone of the plan. Section \ref{sec:synthesis} draws the notions together. Section \ref{sec:proposal} is the proposal, organised in four directions and followed by evaluation, risks and the plan of work.


% Literature review (input by main.tex after the Introduction).
% Sections: foundations; causality as intervention; causality in
% neighbouring fields (HOQC, concurrency); synthesis.

% ---------------------------------------------------------------------
\section{Semantic foundations}\label{sec:foundations}

\subsection{Monads, algebraic theories, and their combinations}

Moggi's observation that notions of computation are monads \cite{MOGGI199155,moggi1991monadicmetalanguage} and Plotkin and Power's refinement that most of them are presented by algebraic theories of operations and equations \cite{plotkin_power_2002,PlotkinPower04CompEffectsOverview} are the foundation of everything below. An algebraic theory determines a monad, and conversely a finitary monad is a Lawvere theory \cite{lawveretheories_monads_HYLAND2007437}; the equational presentation is what makes program equivalence a matter of derivability. Hyland, Plotkin and Power \cite{sums_tensors_HYLAND200670} showed that the standard ways of combining theories have computational meaning: the \emph{sum} adjoins the operations of two theories with no interaction, so that the order of effects from the two theories is observable, while the \emph{tensor} additionally imposes that operations of one theory commute with those of the other, so that their relative order is unobservable. In the premonoidal semantics of effects \cite{PowerRobinson97premonoidal,fuhrmann1999premonoidal}, the morphisms that commute with everything are the \emph{central} ones; centrality is the established algebraic form of ``cannot be observed to happen before or after'', and so is already a primitive statement of causal independence. Handlers \cite{bauer2019algebraicalgebraiceffectshandlers} give the programming counterpart of models; comodels and runners \cite{model_comodel_PLOTKIN2008295,runners_10.1007/978-3-030-44914-8_2} describe the environments in which effectful programs execute. Operations that delimit a scope, such as catching an exception or running a nondeterministic sub-computation once, are not algebraic in Plotkin and Power's sense; Piróg, Schrijvers, Wu and Jaskelioff's \emph{scoped effects} \cite{pirog2018scoped} extend the framework to them, which matters here because normalisation, renormalising a sub-computation after conditioning, is an operation of exactly this kind (\S\ref{sec:prop-graph}).

Ordinary algebraic theories cannot express operations that create and bind \emph{names} or \emph{resources}. Staton's \emph{parameterised algebraic theories} (PATs) \cite{statonlics13,statonpopl15} extend the framework with parameters: operations consume parameters, classically choose a branch, and bind fresh parameters in each branch. Stark's free model of $\pi$-calculus name creation \cite{stark:freamp} is an earlier instance of the same discipline, and Staton used PATs to give complete equational theories of local state \cite{statonlics13} and of qubit quantum computation \cite{statonpopl15}. Kammar, Liell-Cock, Lindley, Matache and Staton \cite{kammar2026threads} have given a complete equational axiomatisation of dynamic threads (fork, wait, stop, atomic actions) as a parameterised algebraic theory, whose closed terms correspond exactly to pomsets; I return to this in \S\ref{sec:concurrency}. PATs are the formalism in which the proposal's algebraic theory of sampling, intervention and normalisation is to be stated (\S\ref{sec:prop-graph}), with intervention labels as parameters.

\subsection{Semantics of probabilistic programs}

Probabilistic programming gives a compositional syntax for statistical models and performs Bayesian inference on them compositionally \cite{vandemeent2021introductionprobabilisticprogramming}. Its denotational semantics is monadic: the Giry monad on measurable spaces \cite{kock72,semantics_ppl_staton2016}, and, for higher-order programs, the probability monad on quasi-Borel spaces \cite{Heunen_2017}. Recursion, which the proposal needs for one reading of cyclic models (\S\ref{sec:prop-expressivity}), is handled by Vákár, Kammar and Staton's domain theory for statistical probabilistic programming \cite{vakar2019domain}, in which quasi-Borel spaces are combined with $\omega$-cpos so that fixpoints and sampling coexist; probabilistic coherence spaces \cite{ehrhard2014pcoh} are the other standard model of probabilistic recursion, and the reference point for full abstraction. Implementations such as LazyPPL \cite{lazyppl}, MonadBayes \cite{monadbayes1,monadbayes2} and GenJAX \cite{becker26genjax} realise the monadic structure directly, LazyPPL through Haskell's laziness, which lets infinite-dimensional models be sampled on demand. For inference at scale, the Python systems Pyro \cite{bingham2018pyro}, on PyTorch, and NumPyro \cite{phan2019composable} and GenJAX \cite{becker26genjax}, on JAX, \emph{vectorise}: independent sub-computations and particles are batched into array operations, so that an inference algorithm is a single traced program; this is the engine the proposal intends to sit on (\S\ref{sec:prop-practice}). Conditioning itself has an algebraic account: Stein and Staton \cite{stein2021exact} treat exact conditioning as an effect with a compositional semantics in Markov categories with conditionals.

The synthetic approach of \emph{Markov categories} \cite{fritz2020synthetic,fritz2019probability} axiomatises the categorical structure needed to talk about probability without measure theory: a symmetric monoidal category with a commutative comonoid on every object, where discarding is natural (the unit is terminal) but copying need not be. In Markov categories with conditionals, disintegration and Bayesian inversion are expressed synthetically \cite{cho2019disintegration,fritz2020synthetic}, and Fritz and Klingler \cite{fritz2023dseparation} prove the soundness of d-separation, and the equivalence of the local and global Markov properties, synthetically; completeness of d-separation in the sense of Geiger, Verma and Pearl is the classical result. The \emph{representable} Markov categories of Fritz, Gonda, Perrone and Rischel \cite{fritz2023representable} are those that arise as Kleisli categories of affine commutative monads, which is the bridge between the synthetic presentation and the monadic one that our languages are interpreted in. This is the \emph{factorisation} face of causality: a Bayesian network is a string diagram in a Markov category, and the conditional independences it entails are read off the diagram.

The second piece of structure the semantics relies on is \emph{grading}. Graded monads \cite{katsumata2014_graded_strong_monad,orchard2020unifying,fujii2016towards,dorsch2019graded} and, more generally, locally graded categories \cite{levy2019locallygraded,sarkis2025string} track an index---an effect annotation, a trace, a resource---through composition. Milius, Pattinson and Schröder \cite{milius2015gradedtrace} use graded monads for trace semantics; Liell-Cock and Staton's $\mathbf{ImP}$ construction \cite{liell2025compositional} uses grading to make imprecise probability compositional. The interventional semantics of \S\ref{sec:intv-wherewestand} is built on an analogous construction in which the grade records which interventions are available.

\emph{What this gives us and what it lacks.} Markov categories give factorisation but no primitive notion of intervention; the diagrammatic causal-model literature (\S\ref{sec:diagrammatic}) adds it by surgery on diagrams. Neither offers a notion of program \emph{context}, which is what contextual equivalence needs; and the monadic semantics in which our languages are interpreted is stated for particular categories rather than synthetically (\S\ref{sec:prop-synthetic}).

% ---------------------------------------------------------------------
\section{Causality as intervention}\label{sec:intv}

\subsection{Structural causal models and Pearl's hierarchy}

Pearl's structural causal models (SCMs) \cite{pearl2009causality,Pearl2018_book_of_why} are systems of equations $X_i = f_i(\mathrm{Pa}_i, U_i)$ in which each endogenous variable is a deterministic function of its parents and of an exogenous noise variable; a prior on the noise makes an SCM a probabilistic model, and the equations make it more than a distribution. The surplus is organised by the \emph{causal hierarchy} \cite{bareinboim:etal20}: rung 1 asks observational questions; rung 2 asks \emph{interventional} questions, answered by replacing the equation for $X$ with a constant (the $\dok$ operator); rung 3 asks \emph{counterfactual} questions about a particular run, answered by abduction, action and prediction, implemented graphically by \emph{twinning}, which duplicates the model into worlds that share the noise \cite{Balke1994}. Both operations are \emph{syntactic transformations of the model}. The Neyman--Rubin potential-outcomes framework \cite{rubin1974} is the other dominant formalism; its correspondence with SCMs is well understood \cite{ibeling2023comparing}. Richardson and Robins' single-world intervention graphs \cite{richardson2013swig} present an intervention by splitting a node into its natural and its fixed value, which is the graphical counterpart of the intervention labels of our languages; they are deliberately single-world, and cross-world quantities such as natural direct effects need further assumptions there. Interventions need not be perfect: the literature on stochastic and soft interventions \cite{dawid2002_stoch_intv,didelez2006_stoch_mech_intv,Eberhardt_Scheines_2007,levis2024stochasticinterventionssensitivityanalysis} replaces an equation by another mechanism rather than a constant.

\paragraph{Graphs, Markov conditions and fine tuning.} The graph of an SCM is derived from its equations: an arrow $X\to Y$ when $f_Y$ reads $X$. The \emph{causal Markov condition} says that the induced distribution factorises along the graph, and \emph{d-separation} \cite{pearl2009causality,spirtes2000causation} reads off the conditional independences that this forces; the \emph{faithfulness} assumption of causal discovery says that there are no others. Faithfulness can fail, and Barrett, Lorenz and Oreshkov \cite{barrett2019quantumcausalmodels} give (their \S2.4) the two-line example that the proposal returns to: $Y = X + \lambda_Y \bmod 2$ with $\lambda_Y$ uniform. Here $X$ is a direct cause of $Y$, since $f_Y$ reads $X$, yet $P(Y\mid \dok(X)) = P(Y)$: no intervention on $X$ is detectable at $Y$, and no signal can be sent from $X$ to $Y$. The independence is \emph{fine tuned}: it holds for one distribution on $\lambda_Y$ and for no other. The distinction between independences forced by the graph and accidental ones is thus the distinction between an intensional datum (the graph) and an extensional one (what interventions reveal), and it is the same distinction as the one between the dataflow of a program and its contextual equivalence class (\S\ref{sec:intv-wherewestand}).

\paragraph{Identifiability and the do-calculus.} When only the graph and observational data are given, the question whether an interventional quantity $P(y\mid\dok(x))$ is determined is \emph{identifiability}. Pearl's \emph{do-calculus} \cite{pearl1995docalculus} is a set of three inference rules, each licensed by a d-separation condition in a graph with some arrows cut, that rewrite interventional expressions into observational ones; Huang and Valtorta \cite{huang2006complete} and Shpitser and Pearl \cite{shpitser2006identification,JMLR:v9:shpitser08a} prove it complete for identification, with an algorithm that either produces a formula or a counterexample. The do-calculus is therefore an \emph{equational theory} of interventional queries indexed by a graph, and the natural object to compare with the program-level equational theories of \S\ref{sec:intv-wherewestand} (\S\ref{sec:prop-graph}). Causal discovery \cite{spirtes2000causation,peters_janzing_schoelkopf_2017_book}, which infers the graph itself from data, remains outside the scope of this thesis.

\paragraph{Cyclic models.} Pearl's models are acyclic. Bongers, Forré, Peters and Mooij \cite{Bongers_2021} extend SCMs to \emph{cyclic} systems of equations and give them meaning through \emph{solution functions}: for a fixed value of the noise, a solution is an assignment satisfying all equations at once, and the model defines a distribution when a solution exists and is unique (or measurably selectable). Forré and Mooij \cite{forre2017markov} replace d-separation by $\sigma$-separation and prove the corresponding Markov properties for cyclic graphs, where the acyclic ones fail. Mooij, Janzing and Schölkopf \cite{mooij2013ode} show where such models come from: a system of ordinary differential equations that converges to an equilibrium induces a cyclic SCM whose solutions are the equilibria, and interventions correspond to clamping a variable. This is one reading of feedback, in which the randomness is drawn once and the cycle is a fixpoint of deterministic equations. The other reading, standard in time-series and dynamic Bayesian networks, unrolls the cycle in time: each iteration draws fresh randomness and the model is an acyclic model over indexed copies of the variables, a stochastic process. The two agree in the deterministic equilibrium case and differ in general. Neither has a programming-language account, in which feedback would be a loop or a recursive definition (\S\ref{sec:prop-expressivity}).

\emph{What this gives us.} The causal structure of an SCM is its dataflow: which variables each equation reads. Interventions probe it, but, by fine tuning, do not determine it. The semantic question is when two models have the same causal structure, and the answer must say which of the two one means.

\subsection{Equivalences of causal models and logical accounts}

Bongers, Forré, Peters and Mooij \cite{Bongers_2021} give the principal account of equivalence for SCMs at all three rungs: two models are observationally, interventionally or counterfactually equivalent when they agree on the corresponding class of queries, with the counterfactual case defined through the two-world twinned model, and the treatment covers cyclic models through their solution functions. At rung 2, several further notions exist. Pearl's causal Bayesian networks are families of interventional distributions compatible with a graph. \emph{Causal spaces} \cite{park2024measuretheoreticaxiomatisationcausality,buchholz2024products} axiomatise rung 2 measure-theoretically: a causal space is a probability space together with a family of \emph{causal kernels}, one for each sub-$\sigma$-algebra, giving the distribution after intervening on the events it measures, subject to consistency axioms; products of causal spaces are defined, and the construction makes no reference to a graph or to mechanisms. \emph{Causal abstraction} \cite{rubenstein2017consistency,beckers2019abstracting,rischel2021abstraction,lorenztull2026abstraction,englberger2025causalabstractionscategoricallyunified} studies maps between models of different granularity, a coarse model being an abstraction of a fine one when interventions on the coarse variables and their images on the fine ones commute with the map; Rischel and Weichwald \cite{rischel2021abstraction} and the later string-diagram line make this a morphism of models in a category. At the level of graphs rather than models, Markov equivalence \cite{verma1990equivalence} is what discovery algorithms can recover. On the logical side, Ibeling and Icard \cite{ibeling2020hierarchy} characterise the three rungs by three probabilistic logics of strictly increasing expressivity whose semantics is defined on SCMs through Pearl's operations, and Icard has proposed probabilistic Turing machines as the computational counterpart of SCMs \cite{icard2017programs,ibeling2019openuniverse}.

\emph{What this gives us and what it lacks.} These are all extensional characterisations of interventional structure, and they are the yardstick against which our contextual equivalences are measured (\S\ref{sec:intv-wherewestand}). None of them is formulated as a contextual equivalence of programs, and none treats higher-order models; the proposal asks how each is related to the graded semantics (\S\ref{sec:prop-relations}).

\subsection{Diagrammatic and categorical causal models}\label{sec:diagrammatic}

Fong's \emph{causal theories} \cite{fong2012causaltheories} present a Bayesian network as a symmetric monoidal category freely generated by its graph, with models as functors into stochastic maps; this is the first place where the graph is the \emph{primitive} and the model is derived from it. Jacobs, Kissinger and Zanasi \cite{jacobs2019causal} present interventions as \emph{string-diagram surgery} in this setting, and introduce \emph{comb disintegration}: the use of combs, processes with a hole, to express the classical causal-inference computations, which is the first appearance of the comb notion of \S\ref{sec:hoqc} in classical causal inference and a diagrammatic form of identification. Lorenz and Tull \cite{lorenz2023causalmodelsstringdiagrams} develop causal models in string diagrams systematically, including open models with inputs, general interventions, counterfactuals, and the identifiability of causal effects treated diagrammatically; a later line presents causal abstraction as natural transformations between such models \cite{lorenztull2026abstraction,englberger2025causalabstractionscategoricallyunified}. Bayesian inference in this style is handled by \emph{normalisation boxes}, which renormalise locally inside a diagram \cite{jacobsszelesstein2025inference}, and Agrawal et al.\ \cite{agrawal2026equational} obtain equational principles for causal Bayesian networks from them. Zhang et al.\ \cite{zhang2026semanticseparation} characterise d-separation as a noninterference-style separation property, which makes explicit a connection implicit in all of this work and in \S\ref{sec:intro}: invariance under interventions is a noninterference property \cite{goguen1982noninterference}, and type systems for noninterference, from the dependency core calculus onward \cite{abadi1999dcc}, are the programming-language way of carrying a dependency graph in the \emph{types} rather than in the semantics.

\emph{What this gives us and what it lacks.} This is the closest existing bridge between the factorisation face and the intervention face of causality, and it is compositional in the functorial sense. Its equational theories are, however, stated per graph, and it has no notion of program context or of higher-order models. It is where the graph-centred direction of the proposal (\S\ref{sec:prop-graph}) starts: the surgery, comb and normalisation-box calculi are the diagrammatic form of the do-calculus, and an algebraic theory in which sampling, intervention and normalisation are all operations is the programming-language form that does not yet exist.

\subsection{Causal probabilistic programming languages}

Causal probabilistic programming aims to make intervention and twinning primitives of an expressive probabilistic language. Libraries such as DoWhy \cite{pywhy,blobaum2024dowhygcm} and the $\dok$ operators of Pyro and PyMC \cite{bingham2018pyro,abrilpla2023pymc} manipulate models as graphs or computation graphs and stop at rung 2 (DoWhy's GCM module reaches rung 3 on fitted mechanisms). Witty et al.\ \cite{witty2019synthesis,witty2023bayesian} treat an intervention as a source transformation that replaces subterms, which is powerful but non-compositional; Brulé's Whittemore \cite{brule2018whittemore} is an early embedded causal DSL.

Four systems reach rung 3 natively. OmegaC \cite{tavares2021language} obtains interventions by a lazy evaluation strategy that retroactively modifies closures, reminiscent of dynamic binding, and counterfactuals by making every program a deterministic function of explicit noise. ChiRho \cite{chirho} implements both as effect handlers over Pyro, with worlds realised by tensor broadcasting along a batch dimension and each intervention doubling the world count, so that its relation to Pearl's twinning is not immediate. MultiVerse \cite{perov2020multiverse} uses trace inspection and re-execution. Kiesel et al.\ \cite{kiesel2023whatif} answer counterfactual queries on ProbLog programs by a transformation proved correct against Pearl's semantics.

\emph{What this gives us and what it lacks.} These systems show that the primitives are implementable at higher order, by non-standard evaluation techniques, and ChiRho shows that they can be vectorised. What they lack is a denotational and equational foundation: a semantics in which intervention is a mathematical operation rather than an evaluation strategy, and an account of when two causal programs are equivalent. That is what my first-year work supplies; what it does not yet supply is a semantics \emph{of} these systems, or an implementation that runs at their scale (\S\ref{sec:prop-practice}, \S\ref{sec:prop-relations}).

\subsection{Where I stand}\label{sec:intv-wherewestand}

Our paper, to appear at POPL 2027 \cite{wang2026causalppl}, with Dario Stein (joint first author), Michael Lee, Eli Bingham, John Feser, Ohad Kammar, Sam Staton and Jeremy Yallop (an early version appeared at the LAFI workshop \cite{lafi-wang}; artifact \cite{ARTIFACT_wang_2026_23207221}), reconstructs Pearl's hierarchy as a hierarchy of three first-order probabilistic languages. $\Lk{1}$ is a standard first-order probabilistic language with conditioning by a final score. $\Lk{2}$ exposes hooks in the dataflow as typed free variables, the \emph{intervention labels}, read through $\mathsf{applyIntv}$; an intervention $\dok(\ell\mapsto v)\{t\}$ is substitution of $v$ for $\ell$. $\Lk{3}$ annotates each binding as exogenous or endogenous, and \emph{twinning} $\mathsf{twin}_n$ is a vectorisation transformation from $\Lk{3}$ to $\Lk{2}$ that duplicates the endogenous bindings into $n$ worlds sharing the exogenous ones. Program equivalence at rung $k$ is \emph{contextual} with respect to rung-1 denotational equality, in Milner's sense: two programs are $\equiv_k$ when they are denotationally equal in every rung-$k$ context, and varying the admissible interventions and world counts gives a spectrum of equivalences. For representable acyclic SCMs compiled into $\Lk{3}$, the observational, interventional and counterfactual equivalences of Bongers et al.\ coincide with the equivalences $\equiv_3^{\mathrm{obs}}$, $\equiv_3^{\mathrm{intv}}$ and $\equiv_3^{\text{2-world}}$ of the compiled programs when the model is concrete, the forward implication holding in general; the full $\equiv_3$ quantifies over all $n$-world twinnings, and whether it is strictly finer than $\equiv_3^{\text{2-world}}$ is open. The denotational semantics interprets interventions by \emph{grading}, in a construction reminiscent of $\mathbf{ImP}$ \cite{liell2025compositional}, and twinning by world-polymorphic vectorisation. It is adequate for the contextual equivalences unconditionally and fully abstract when the underlying model is \emph{concrete}, meaning that values correspond to global elements and global elements separate morphisms. The construction is parametric in an ``$\Lk{1}$-model'': a distributive category with a commutative, affine probability monad and enough structure to interpret the base language, with cartesian closure added for the higher-order language; $\mathbf{Set}$, standard Borel spaces and quasi-Borel spaces are the instances, and all three are concrete. The denotation then motivates a higher-order language $\LHO$ with cross-world interventions, into which $\Lk{3}$ compiles in a semantics-preserving way, and which has a definitional Haskell implementation. \fillin{one sentence on which parts are yours: e.g.\ the contextual-equivalence framework and the full-abstraction proofs}

Two consequences frame the proposal. First, $\equiv_2$ is the extensional notion: it identifies exactly the programs that no interventional observer can tell apart. The fine-tuned model of Barrett, Lorenz and Oreshkov is a one-line program, $\mathtt{let}\ \lambda = \mathtt{bernoulli}\ 0.5\ \mathtt{in}\ \mathtt{return}\ (\mathtt{xor}\ t'\ \lambda)$ with $t' = \mathtt{applyIntv}\ \ell\ t$, and it is $\equiv_2$ to the program that never reads $t'$; the dataflow, which is the graph, is finer than $\equiv_2$, and the full-abstraction theorem is a theorem about the extensional notion only. Second, the limitations relevant to the proposal are these: the languages are acyclic; the relation between the contextual equational theory and the graph-indexed equational theories of \S\ref{sec:diagrammatic}, the do-calculus included, is not worked out; the categorical assumptions are a list of properties of particular categories rather than a synthetic axiomatisation; the first-order semantics is over normalised probability, while $\LHO$ widens to the unnormalised measure monad to internalise conditioning, so that programs with observations denote unnormalised, weighted maps; and the models are measure-theoretic rather than finite-dimensional, which matters for the signalling face (\S\ref{sec:synthesis}).

% ---------------------------------------------------------------------
\section{Causality in neighbouring fields}\label{sec:order}

The two fields of this section are positioned, not surveyed: my results are not in them, I know them at the level of their main definitions and results, and their detailed study is a milestone of the plan (\S\ref{sec:plan}). They are included because each has a notion of causality that the proposal's fourth direction (\S\ref{sec:prop-other}) wants to relate to the interventional one. The one item of that direction that is a theorem rather than a position, the classical bridge D2, needs only Kissinger and Uijlen's first-order and one-way-signalling types and their classical instance over $\Mat$, which I do know.

\subsection{Higher-order quantum processes and signalling}\label{sec:hoqc}

In the standard semantics of quantum computation \cite{QPL_SELINGER_2004,PQP_Coecke_Kissinger_2017}, types denote finite-dimensional spaces of density matrices and programs denote completely positive maps; the deterministic processes, those that preserve the discarding map, are the trace-preserving ones, and the category $\CP$ of all completely positive maps is compact closed, which the subcategory of deterministic ones is not. Higher-order quantum computation studies transformations of such processes rather than of states. Chiribella, D'Ariano and Perinotti's \emph{quantum combs} \cite{quantum_circuit_architecture_Chiribella_2008,chiribella2009networks} are processes with holes in a fixed causal order; their probabilistic counterparts, \emph{testers}, and Gutoski and Watrous's \emph{strategies} \cite{gutoski2007games} describe multi-round interactions with sub-normalised outcomes, and the general theory of supermaps forms a hierarchy of types \cite{bisio2019hoqt}. Oreshkov, Costa and Brukner's \emph{process matrices} \cite{oreshkov2012noncausal} are the most general bipartite process operators compatible with quantum mechanics holding locally in each party's laboratory; the multipartite theory and the notion of \emph{causal separability}, in which the order of later parties may depend classically on earlier outcomes, are due to Araújo et al.\ and Oreshkov and Giarmatzi \cite{araujo2015witnessing,oreshkov2016separable}. Some process matrices are causally non-separable; the \emph{quantum switch} \cite{chiribella2013switch}, which applies two operations in an order controlled by a qubit in superposition, is the standard example \cite{araujo2015witnessing}, and which indefinite-order processes can be built as circuits with classical or quantum control of order is the subject of Wechs et al.\ \cite{wechs2021qcqc}.

Categorically, Kissinger and Uijlen's $\CausC$ construction \cite{kissinger2019categorical} builds, over a \emph{precausal} category---compact closed, with a chosen discarding structure and further axioms, the examples being $\CP$ on finite-dimensional systems and $\Mat$, finite matrices of non-negative reals---a $*$-autonomous category whose types are signalling constraints: first-order processes, non-signalling processes, one-way signalling processes, and so on. The classical example matters for this thesis: every construction of the quantum theory goes through for $\Mat$, where first-order processes are stochastic maps, so that combs, testers and signalling types have a purely classical meaning. The dual of the non-signalling type contains all bipartite process matrices, causally separable or not; separability is a property inside that type, not a type, and Simmons shows, as \emph{affine universality} of definite orders, that every process matrix is an affine combination of definite-order processes \cite{simmons2024thesis}. Simmons's thesis \cite{simmons2024thesis} is the fullest development of this line. It rebuilds the construction over \emph{additive} precausal categories, which gives $\CausC$ products, coproducts and probabilistic choice as type operators; it adds a non-commutative sequential operator for one-way signalling, under which $\CausC$ models Guglielmi's BV-logic \cite{simmons2022bv}; it defines $n$-ary \emph{graph types} encoding compatibility with an arbitrary DAG of signalling relations, which is the one existing place where a causal graph is a \emph{type}; it surveys candidate ways of admitting sub-normalised, partial processes into the construction and where each falls short; and it introduces \emph{causal logic}, pomset logic with directed axiom links, whose proof nets characterise exactly the causally consistent wirings of black boxes \cite{simmons2024thesis,simmons2024completelogic}. Simmons's finite-causal-basis axiom excludes infinite-dimensional and measure-theoretic bases, which is one of the two obstructions discussed in \S\ref{sec:synthesis}.

\emph{Quantum causal models} transfer Pearl's framework to quantum processes: Costa and Shrapnel \cite{costa2016quantumcausalmodelling} gave the first process-matrix-based definition, Allen et al.\ \cite{allen2017quantumcommoncauses} the common-cause line, and Barrett, Lorenz and Oreshkov \cite{barrett2019quantumcausalmodels} the general framework, in which a node is an input--output pair, the process operator factorises as a product of commuting operators (a quantum Markov condition), interventions are local instruments at nodes, d-separation in the graph is sound and complete for a quantum analogue of conditional independence, and the three rules of the do-calculus have quantum generalisations. Three features of their framework (their \S\S2.4, 4.1 and 4.5) are relevant here. Causal structure is defined for a \emph{unitary} realisation by a no-influence condition between inputs and outputs, and the DAG of a model is derived from it; the DAG is in general a strict subgraph of the \emph{connectivity} DAG of the circuit, since a unitary may have no-influence relations that its wiring does not show; and, as in the classical case, a direct cause need not permit signalling in the marginal process, the signalling being washed out by a fine-tuned state on an unobserved system. There are thus three strictly nested notions, wiring, causal structure and signalling, of which the middle one is the primitive. The cyclic extension \cite{barrett2021cyclic} gives a cyclic-causal-model account of some processes without definite order, including the switch, which is the quantum form of the observation that feedback and indefiniteness of order are related.

\emph{Position.} Two things follow for this thesis. First, signalling types are the extensional face of causal structure in the same way that interventional equivalence is the extensional face of dataflow, and the fine-tuning phenomenon is the same in both: a signalling constraint says that a local intervention at one node is undetectable at another, which is the interventional notion stated for processes with slots. Second, quantum causal models are the one existing place where factorisation, intervention and signalling are all present, and they are second-order processes with nodes as slots, exactly as a program with intervention labels is a process with slots; that is why they are a reading milestone and the starting point of \S\ref{sec:prop-hoqc}.

\subsection{Order in concurrency theory}\label{sec:concurrency}

Concurrency theory has long modelled causal order explicitly. Lamport's happens-before relation \cite{lamport1978time} is the partial order an observer of a distributed system can reconstruct. Mazurkiewicz traces \cite{diekert1995traces} are words modulo the commutation of independent letters, the cheapest partial-order model. Pratt's pomsets \cite{pratt1986pomsets}---partially ordered multisets of events, up to isomorphism, with series and parallel composition---generalise words to partially ordered words, and Gischer \cite{gischer1988pomsets} axiomatised the equational theory of series-parallel pomsets. Winskel's event structures \cite{winskel1986eventstructures,nielsen1981petri,winskel1995models} add a conflict relation to the causal order, so that a single structure records both what must precede what and what excludes what, and give a domain-theoretic account of nondeterministic concurrent computation. Concurrent Kleene algebra \cite{hoare2011cka} adds a commutative parallel operator to Kleene algebra, related to sequential composition by an exchange law; Kappé, Brunet, Silva and Zanasi \cite{kappe2018cka} prove that series-parallel pomset languages closed under Gischer's subsumption order are its free model, building on Laurence and Struth's completeness theorem for bi-Kleene algebra \cite{laurence2014completeness}. Within the algebraic-effects tradition, Dvir, Kammar and Lahav \cite{dvir2022sharedstate} give a monadic semantics of shared-state concurrency from an equational theory, and Kammar et al.\ \cite{kammar2026threads} axiomatise dynamic threads as a parameterised algebraic theory whose closed terms are exactly pomsets, with a completeness theorem. Clairambault, de Visme and Winskel \cite{clairambault2019quantumgames} give an event-structure game semantics for quantum programs.

\emph{Position.} In these models the causal graph is \emph{explicit}: a pomset or event structure is a graph of events, and the semantics is the graph. This is the opposite extreme from the interventional line, where the graph is implicit in the dataflow and only its extensional shadow is characterised. Whether the interventional notion has a counterpart in these models, intervening on an event, and whether the signalling types of \S\ref{sec:hoqc} are the extensional reading of a pomset, are the open-ended questions of \S\ref{sec:prop-concurrency}.

% ---------------------------------------------------------------------
\section{Synthesis: three notions of causality and the gap}\label{sec:synthesis}

Table \ref{tab:synthesis} summarises the review. Each column is one extensional notion of causality; each row asks the same question of it.

\begin{table}[ht]
\centering\small
\begin{tabularx}{\textwidth}{>{\raggedright\arraybackslash}p{2.6cm}>{\raggedright\arraybackslash}X>{\raggedright\arraybackslash}X>{\raggedright\arraybackslash}X}
\toprule
 & \textbf{Intervention} & \textbf{Factorisation} & \textbf{Signalling}\\
\midrule
Intensional datum & dataflow of a program / SCM & DAG & wiring of a circuit; comb order; pomset\\
Extensional characterisation & contextual equivalence under interventions & conditional independence / d-separation & signalling types ($\CausC$); causal (non-)separability\\
Gap between the two & fine tuning: dataflow finer than $\equiv_2$ & faithfulness: graph finer than the independences & fine tuning: causal structure finer than signalling \cite{barrett2019quantumcausalmodels}\\
Correspondence theorem & full abstraction for concrete models (\S\ref{sec:intv-wherewestand}) & soundness of d-separation, synthetically \cite{fritz2023dseparation}; completeness classically & types of $\CausC$ characterise signalling; no programming-language counterpart\\
Equational theory & contextual equivalences $\equiv_k$ & do-calculus; normalisation-box equations & BV and causal logic\\
Semantics of programs & first-order and higher-order, acyclic & Markov categories; interventions by surgery & none for signalling; pomsets for concurrent programs\\
Next hard case & cyclic SCMs & cyclic and quantum causal models & causally non-separable processes; conflict\\
\bottomrule
\end{tabularx}
\caption{Three notions of causality and where the first-year work stands.}\label{tab:synthesis}
\end{table}

Four remarks. First, the asymmetry: for intervention the correspondence between intensional and extensional structure is a theorem, and the intensional datum is a program; for signalling the extensional characterisation exists, as a type system, but there is no programming language whose syntax it characterises; for concurrency the intensional datum is explicit and there is no interventional reading of it. Second, the row ``gap between the two'' is the same phenomenon in all three columns, and it is what motivates the graph-centred direction of the proposal: in each column the intensional datum is strictly finer than what any observer can recover, and the literature's response in each is to take the graph as primitive, as a Markov condition, a no-influence condition, or a type. Third, non-DAG structure is the next hard case in every column, but it is not the same case: feedback in cyclic SCMs, causal non-separability in quantum processes, and conflict in event structures are three different departures from a fixed order, and a unifying account must distinguish them. Fourth, and this is the gap: there is no framework in which the three columns are structure on one category. Markov categories cover factorisation and, through surgery, part of intervention; $\CausC$ covers signalling for quantum and finite classical processes; quantum causal models relate all three, with a compositional presentation in Lorenz and Tull's string diagrams, but are not formulated as a semantics of programs. The graded semantics of \S\ref{sec:intv-wherewestand} covers the first column; it is a family of morphisms indexed by an intensional datum, the grade, with a reindexing action, which is the common structure a generalisation would start from.

Two obstructions should be displayed rather than hidden. The first is \emph{finite dimension}: the $\CausC$ construction needs a compact closed base with invertible dimensions, which $\Mat$ and finite-dimensional $\CP$ provide, but measurable spaces and quasi-Borel spaces do not. Any bridge from the interventional semantics to signalling types is therefore stated on the finite fragment. The second is \emph{sub-normalisation}: conditioning in $\LHO$ denotes unnormalised, weighted maps, which do not land in a first-order $\CausC$ type as it stands; testers and strategies \cite{chiribella2009networks,gutoski2007games} accommodate sub-normalised families only, so a renormalisation is needed even to reach that setting, Simmons \cite{simmons2024thesis} surveys the ways of admitting partial processes into $\CausC$ and finds each wanting in some respect, and normalisation boxes \cite{jacobsszelesstein2025inference} are the diagrammatic tool for renormalising after conditioning.


% =====================================================================
\section{Project Proposal}\label{sec:proposal}
% =====================================================================
% Structure follows main.tex: four directions, then evaluation/risks
% and the plan. The handbook's questions (i)--(vi) are answered in the
% opening subsection and in the two closing ones; each item in the four
% directions carries a risk rating and a "done" criterion.

\subsection{Topic, goal, state of the art, and method}

\paragraph{Research topic and fundamental challenges.}
The thesis is on the \emph{semantics of causal probabilistic programming}: a programming-language account of causal models in which the interventional face of causal structure is defined by a contextual equivalence of programs, interventions and counterfactuals are mathematical operations on denotations rather than evaluation strategies, and the account is executable. The finished work of \S\ref{sec:intv-wherewestand} is the base. The thesis builds outward from it in four directions, which are the four subsections below: making the language practical and more expressive; understanding the categorical constructions it rests on; taking the causal graph, rather than the model, as the primitive; and relating the interventional notion of causality to the notions used in concurrency and in higher-order quantum processes. The fundamental challenges are, respectively: (i) cycles, since nothing in the graded semantics handles feedback, and feedback interacts with interventions in ways acyclic models never exhibit; (ii) the semantics is parametric in a list of categorical assumptions rather than in a synthetic axiomatisation, and its relation to the other extensional accounts, causal spaces, causal abstraction and the lazy semantics of OmegaC, is unproven; (iii) the extensional notion $\equiv_2$ is blind to fine-tuned dependence, so a graph-centred account must add intensional structure to the types or to the semantics, and then show how the do-calculus lives there; (iv) the signalling notion of causality is formulated in compact closed categories over finite-dimensional processes, and the interventional one in cartesian closed categories over measure-theoretic ones, so a relation between them must cross that divide.

\paragraph{Goal.}
I want a mathematical account of what it means for one part of a program to cause another, in the sense that changing it changes the other. It should work for programs that pass procedures around and for programs that loop back on themselves; it should run on the software that statisticians actually use; and it should explain why the same word, cause, is used for the ability to send a message between parts of a quantum process and for the order of events in a concurrent one.

\paragraph{State of the art and why it is hard.}
The existing notions of equivalence for causal models, Bongers et al.'s hierarchy, causal spaces and causal abstraction, say when two models have the same interventional behaviour, but they are not equivalences of programs. Causal probabilistic programming languages, OmegaC, ChiRho and MultiVerse, implement interventions and counterfactuals at higher order by evaluation strategies without an account of what they compute, and ChiRho alone runs at scale. The diagrammatic line has equational principles for causal Bayesian networks and a diagrammatic do-calculus, stated per graph, with no notion of program. Cyclic models have a semantics only at the model level, by solution functions or by unrolling in time. On the other side, higher-order quantum processes have a type-theoretic account of signalling, and concurrency has explicit causal graphs, with no account of intervention in either. The problem is hard because each extension breaks something the base relies on: cycles break the dataflow reading; a synthetic setting removes the global elements that full abstraction used; a graph-centred account must reinstate exactly what contextual equivalence abstracts away; and the signalling setting lacks the function spaces that the higher-order semantics lives in.

\paragraph{Method and what is new.}
The method is the one that produced the base result. (1) Fix a small language whose \emph{syntax} carries the causal structure. (2) Define an \emph{extensional} equivalence from what an observer may do. (3) Build a \emph{denotational} model in which the available observations are a grade. (4) Prove a correspondence theorem between (2) and (3). What is new is that the interventional face of causal structure is \emph{defined} by a contextual equivalence of programs rather than by a property of a graph, which is what lets the account survive higher order, cycles and abstraction, and that every semantic definition is also an implementation. Where the direction is open-ended (\S\ref{sec:prop-other}), I say so and commit to reading and a position rather than to a theorem.

\paragraph{Running example.} In $\Lk{2}$, interventions are exposed as typed free variables $\ell$ read through $\mathsf{applyIntv}$, which returns the label's value when it is set and the program's own value otherwise:
\begin{lstlisting}
let t  = bernoulli 0.5 in              -- treatment
let t' = applyIntv $\ell$ t in             -- intervenable copy of t
let m  = bernoulli (f t') in           -- mediator
let y  = bernoulli (g t' m) in         -- outcome
let z  = bernoulli 0.3 in              -- unrelated variable
return (y, z)
\end{lstlisting}
Under $\dok(\ell\mapsto\mathsf{Just}\ \mathsf{true})\{-\}$ the variable $t'$ denotes the constant $\mathsf{true}$, and the total effect of $t$ on $y$ is the difference between the distributions of $y$ under the two perfect interventions. The natural direct effect, a cross-world quantity, is expressed in $\LHO$ by twinning the program and intervening on the mediator of one world with a value read from the other. The variable $z$ has the same distribution in every interventional context, and $\equiv_2$ makes this an equation between programs. Each direction below modifies this example: adding feedback ($m$ reads $y$), coarsening ($t$ and $m$ into one variable), fine tuning ($y$ reads $t'$ but no intervention shows it), and reading $\ell$ as the slot of a higher-order process.

% ---------------------------------------------------------------------
\subsection{Developing the field of causal probabilistic programming}

\subsubsection{Putting the language into practice}\label{sec:prop-practice}

The existing implementation is a definitional library in Haskell over MonadBayes, a shallow embedding that handles small examples only. Realistic causal inference needs vectorised inference, and the candidate host is Python, on top of NumPyro or JAX directly (\S\ref{sec:foundations}), which gives an inference engine that is efficient by construction and the user base that causal machine learning actually has. The point is not a port but a \emph{library with guarantees}: the operations of $\LHO$, intervention labels, $\dok$, twinning and cross-world interventions, implemented as program transformations whose correctness is the denotational semantics, with as much of the typing discipline (linearity of labels, world indices) as the host allows, enforced statically where Python's type hints permit and dynamically otherwise.
\begin{itemize}[nosep]
\item[A1] \textit{(low risk; minimum viable thesis)} The library: $\LHO$'s constructs over a vectorised host, with twinning implemented as a batch dimension realising $\LHO$'s world-indexed types, ChiRho's tensor broadcasting being the point of comparison. \emph{Done:} the paper's examples at realistic scale, and the semantics stated as the specification of each construct. \emph{Why it may succeed:} the semantics already is a vectorisation, so the batch dimension is its direct image.
\item[A2] \textit{(medium risk)} Evaluation against ChiRho on a benchmark set \fillin{to be agreed with supervisor}. \emph{Done:} the set fixed; accuracy and wall-clock time measured on it; a table of which $\LHO$ guarantees survive the host, statically, dynamically, or not at all.
\end{itemize}

\subsubsection{Expressivity: cycles, feedback and recursion}\label{sec:prop-expressivity}

Replace the mediator's and outcome's equations in the running example by ones that read each other, with the noise drawn outside the loop: $\mathtt{let}\ u_m = \mathtt{uniform}\ \mathtt{in}\ \mathtt{let}\ u_y = \mathtt{uniform}\ \mathtt{in}\ \mathtt{fix}\ (m,y).\ (u_m < f\ t'\ y,\ u_y < g\ t'\ m)$, so that $m$ and $y$ depend on each other. No acyclic program denotes this, and for Boolean $m, y$ the system has a unique solution only for some $f, g$ and some values of the noise, which is the first thing a semantics must say. The review (\S\ref{sec:intv}) distinguishes two readings of such a cycle, and the proposal is to give both a programming-language semantics and to relate them.

\emph{Reading 1: equations.} The randomness is drawn once and the cycle is a fixpoint of deterministic equations; this is Bongers et al.'s solution-function semantics and the equilibrium reading of Mooij et al. In a language with a fixpoint construct $\mathtt{fix}\ (m,y).\ (\dots)$ over an explicit noise variable, the model is a program again. This reading is easier precisely because the model is \emph{equational}: an SCM in $\Lk{3}$ already separates exogenous from endogenous bindings, so the fixpoint is taken over the endogenous ones with the noise fixed; the body of the fixpoint must be deterministic in the noise, as in the example, since $\Lk{3}$ allows endogenous bindings that sample, and those are kept out of the loop in the first instance. Intervening on a fixpoint-bound variable is cutting its equation and re-solving, which requires unique solvability after every intervention, Bongers et al.'s \emph{simple} SCMs, and the program-level definition must agree with theirs where both are defined. The same construction, closing an output onto an input, is a \emph{feedback loop} in a compact closed category, which is how processes without definite order arise in the quantum setting and how Barrett, Lorenz and Oreshkov's cyclic quantum causal models are built (\S\ref{sec:hoqc}); the relation is pursued in \S\ref{sec:prop-hoqc}.

\emph{Reading 2: recursion.} The randomness is resampled at each pass through the cycle, so that the model is a stochastic process, and the cycle is a loop or a recursive definition in the probabilistic language, with Vákár, Kammar and Staton's domain theory as the natural semantics and laziness, as in LazyPPL and in OmegaC's retroactive interventions, as the natural implementation. Recursive types, streams in particular, and mutual recursion are the constructs to examine: a cyclic model over streams of variables is an acyclic model unrolled in time, which is the dynamic-Bayesian-network reading. The question for interventions is which pass one intervenes on, and whether an intervention on a stream is itself a stream of interventions.
\begin{itemize}[nosep]
\item[A3] \textit{(medium risk; minimum viable thesis)} Reading 1 on finite simple models: a fixpoint construct over $\Lk{3}$ with noise-deterministic bodies, its denotation as the solution function, interventions by cutting and re-solving, the grade extended, adequacy proved, and full abstraction proved over $\mathbf{Set}$-models, where finiteness makes concreteness free. \emph{Done:} the extension, the two theorems, and the feedback example reproduced for the $f, g$ that make it simple. \emph{Failure:} full abstraction fails for the cyclic language, which would show that cycles need a new observable. \emph{Why it may succeed:} the solution-function semantics is already a fixpoint; the question is which construct denotes it.
\item[A4] \textit{(high risk; contingent on A3)} Reading 2, and the relation between the readings: a recursive $\Lk{2}$ over $\omega$-quasi-Borel spaces with intervention labels on streams, with denotational $\Lk{1}$-equality in that model as the observable. \emph{Done:} the agreement theorem, that the two readings coincide on deterministic equilibria, stated and proved, with a characterisation of where they differ; full abstraction for the recursive language, or a precise obstruction.
\end{itemize}

% ---------------------------------------------------------------------
\subsection{Better understanding our categorical constructions}

\subsubsection{Relation to other constructions}\label{sec:prop-relations}

The graded interventional construction and its higher-order internalisation are new objects, and I do not yet know how expressive they are or how they relate to the other extensional accounts in \S\ref{sec:intv}. Three relations are concrete enough to state.

\emph{OmegaC.} My hypothesis is that OmegaC's lazy semantics is the interventional monad applied \emph{pointwise to each value}: where my semantics threads the intervention environment through a program's dataflow once, a lazy semantics attaches it to every thunk, so that a retroactive modification of a closure is a reindexing of the grade at the point of forcing. If so, OmegaC's core has a denotational model in the same construction, its evaluation strategy computes $\dok$, and the two semantics are related by a call-by-need translation.
\begin{itemize}[nosep]
\item[B1] \textit{(medium risk; the OmegaC half in the minimum viable thesis)} A formal operational semantics for OmegaC's core, a denotational model of it in the graded construction, and the correctness of its lazy evaluation strategy with respect to $\dok$. For ChiRho, whose worlds are made by tensor broadcasting with each intervention doubling the world count, the expected result is a precise \emph{mismatch} theorem: which values ChiRho shares across worlds that $\LHO$ twinning does not, and conversely. \emph{Done:} the translations, their correctness or mismatch statements, and the papers' examples reproduced. \emph{Why it may succeed:} OmegaC is by construction a deterministic function of explicit noise, which is what $\Lk{3}$'s exogenous/endogenous split formalises; ChiRho's broadcasting rule is a computable function of the execution trace, which is all a mismatch theorem needs.
\end{itemize}
\emph{Causal spaces.} A causal space is a probability space with a causal kernel for every sub-$\sigma$-algebra. The denotation of a closed $\Lk{2}$ program at the full grade is a distribution together with a reindexing action for every set of labels, which looks like a causal space on the $\sigma$-algebra generated by the labelled variables.
\begin{itemize}[nosep]
\item[B2] \textit{(medium risk)} A functor from $\Lk{2}$ denotations to causal spaces, and a characterisation of its image. \emph{Done:} the functor, with the consistency axioms of causal spaces derived from the laws of the grade, or a counterexample showing which axiom the grade does not force.
\end{itemize}
\emph{Causal abstraction.} Coarsening $t$ and $m$ of the running example into a single variable is an abstraction. At the program level it is a transformation between $\Lk{2}$ programs that commutes with interventions on the coarse labels.
\begin{itemize}[nosep]
\item[B3] \textit{(medium risk)} Causal abstraction as a program transformation, specialising to the model-level definitions \cite{rischel2021abstraction,lorenztull2026abstraction} on compiled SCMs, which asks the graded semantics to be functorial in the model and not only in the program. \emph{Done:} the definition, the specialisation theorem, and the worked example.
\end{itemize}

\subsubsection{A more synthetic approach?}\label{sec:prop-synthetic}

The semantics is stated for an ``$\Lk{1}$-model'': a distributive category with a commutative affine probability monad and the structure needed for the base language, with cartesian closure for the higher-order part, instantiated at $\mathbf{Set}$, standard Borel spaces and quasi-Borel spaces. This is a shopping list. The synthetic alternative is to work in a Markov category, where the \emph{factorisation} face of causality is already native and interventions are diagram surgery: the intervention label $\ell$ becomes an open input wire, $\mathsf{applyIntv}$ a diagram that merges it with the program's own value, and the grade the set of wires left open. The representable Markov categories of Fritz et al.\ \cite{fritz2023representable}, Kleisli categories of affine monads, are where the two presentations meet. What may be lost is full abstraction, which used global elements to separate morphisms; in a Markov category the substitute is a well-pointedness or enough-points axiom, and finding the weakest one that suffices is the content.
\begin{itemize}[nosep]
\item[B4] \textit{(medium risk; minimum viable thesis)} The graded interventional construction over a Markov category, agreeing with the current one on representable Markov categories, with the axioms needed for adequacy and for full abstraction identified. \emph{Done:} the construction, the agreement theorem, and the axiom list, which replaces the shopping list and will include finite coproducts with distributivity, since Maybe-typed labels and $\mathsf{applyIntv}$ need them. \emph{Why it may succeed:} the ingredients the proofs use, copying, discarding and value separation, are each a named axiom in the Markov-category literature. \emph{Failure:} if no axiom weaker than concreteness suffices for full abstraction, that is the result.
\end{itemize}

% ---------------------------------------------------------------------
\subsection{Understanding causality using a graph-centred approach}\label{sec:prop-graph}

\paragraph{The graph as primitive.} Barrett, Lorenz and Oreshkov's example (\S\ref{sec:intv}) in the running language: $y = \mathtt{xor}\ t'\ \lambda$ with $\lambda$ uniform. The program reads $t'$, so the graph has an arrow from $t$ to $y$; no intervention on $\ell$ changes the distribution of the output, so the program is $\equiv_2$ to one that never reads $t'$. Contextual equivalence is a \emph{signalling} notion, with perfect interventions as the signals, and it is blind to fine tuning. Separate the signature into the built-ins ($\mathtt{if}$, $\mathsf{applyIntv}$, arithmetic) and the \emph{model symbols}, the mechanisms $f, g$ and the priors, whose interpretation the base semantics leaves open; write the example with the prior of $\lambda$ as a model symbol $p$, realised at $\mathtt{bernoulli}\ 0.5$. A \emph{realisation} is an interpretation of the model symbols. An arrow of the dataflow graph is \emph{faithful} when some realisation lets an intervention on its source be detected at its target, and \emph{unfaithful} at a given realisation when the dependence is genuine but invisible there; the arrow from $t$ to $y$ is faithful, since any realisation of $p$ other than the uniform one reveals it, and the program at the uniform realisation is fine tuned. The interventional base result characterises invariance at a fixed realisation only. A graph-centred account takes the graph as the datum and asks what structure in the language or the semantics carries it. Two options are available, and the proposal tries both on the first-order language.

\emph{Graph in the semantics.} Fong's causal theories and the string-diagram models of Jacobs et al.\ and Lorenz and Tull take the graph as the generating data of a category and the model as a functor; a program's denotation would then be a functor out of the free Markov category on its dataflow graph, and two programs would be equivalent when their dataflow graphs coincide and the functors agree, which identifies strictly less than $\equiv_2$ and is, deliberately, close to syntactic equality of mechanisms.

\emph{Graph in the types.} A dependency type system, in the tradition of the dependency core calculus \cite{abadi1999dcc}, records in the type of each variable which labels it depends on; Simmons's graph types for higher-order processes (\S\ref{sec:hoqc}) are the quantum analogue, with a DAG of signalling relations as a type. A graph-typed $\Lk{2}$ would have a soundness theorem, that a typed edge implies signalling under some realisation, and fine tuning would be exactly the gap between the typed graph and $\equiv_2$ at a fixed realisation.
\begin{itemize}[nosep]
\item[C1] \textit{(medium risk)} A graph-typed first-order language with both semantics, the functorial one and the graded one, related by a theorem: there is an edge from $X$ to $Y$ in the typed graph iff some realisation of the program exhibits signalling from $X$ to $Y$. \emph{Done:} the type system, the two semantics, the relating theorem, and the fine-tuning example as the separating case. \emph{Why it may succeed:} the typed graph is a dependency analysis, and dependency analyses are sound for noninterference by standard arguments \cite{abadi1999dcc}; the content is the converse.
\end{itemize}

\paragraph{Identifiability as an equational theory.} The do-calculus is an equational theory of interventional queries indexed by a graph, and it is interesting only when one has a causal structure and nothing more; its diagrammatic forms are comb disintegration, Lorenz and Tull's identifiability, and the normalisation-box equations (\S\ref{sec:diagrammatic}). None of these is an algebraic theory of \emph{operations}. The proposal is a parameterised algebraic theory with three kinds of operation: sampling, as in Staton's theory of probability; intervention, with labels as parameters; and normalisation, as a \emph{scoped} operation in the sense of Piróg et al.\ \cite{pirog2018scoped}, delimiting the sub-computation to be renormalised after conditioning, in the manner of Stein and Staton's exact conditioning \cite{stein2021exact}. Everything is then in one framework: for terms annotated with a given DAG, the graph being a parameter of the theory as in Fong's causal theories, the Markov conditions and d-separation become equations of the theory, and the three rules of the do-calculus become candidate derived rules; C1's type system is what later \emph{infers} the annotation, so C2 does not depend on C1. The prerequisites, the Markov conditions, d-separation and the calculus itself, are covered by the first reading milestone M1.
\begin{itemize}[nosep]
\item[C2] \textit{(medium risk; minimum viable thesis)} The theory, parameterised by a DAG, with a sound interpretation in Markov categories with conditionals, and Agrawal et al.'s equational principles for causal Bayesian networks, obtained from normalisation boxes, derived in it. \emph{Done:} the theory, the soundness theorem, and the derivations. \emph{Failure:} a principle of Agrawal et al.\ that is not derivable, which would locate a missing equation. \emph{Why it may succeed:} normalisation as a scoped operation already has an algebraic treatment \cite{pirog2018scoped}, and exact conditioning already has a sound one \cite{stein2021exact}; the theory combines them with surgery.
\item[C3] \textit{(medium--high risk)} Admissibility of the three do-calculus rules for DAG-annotated terms, each rule's d-separation side condition becoming a condition on the annotation. \emph{Done:} the admissibility theorem, or a precise statement of which rule fails and why. \emph{Why it may succeed:} the soundness of each rule should follow from Fritz and Klingler's synthetic d-separation together with a surgery lemma for the intervention operation, which C2 must prove in any case; the open question is whether the theory's equations suffice to derive the rules, which is why the rating is medium--high.
\end{itemize}

% ---------------------------------------------------------------------
\subsection{Understanding the relationship with `causality' as used in other fields}\label{sec:prop-other}

This direction is deliberately open-ended. Its goal is a categorical account that can speak of causality in the interventional, the concurrent and the higher-order quantum sense at once, and the honest statement of its status is that I have a semantics on one side and an analogy on the other. I commit here to reading milestones, to one classical theorem, and to a position; anything beyond that is a stretch goal.

\subsubsection{Concurrency}\label{sec:prop-concurrency}

Concurrency has the causal graph \emph{explicitly}: a pomset is a graph of events, and event structures add conflict, so that causal inconsistency, an event that both requires and excludes another, is representable. The interventional line has the opposite profile, an implicit graph with an extensional characterisation, and the two have not been put side by side. Two questions. First, what is an intervention on an event, and is the order of a pomset the least order compatible with every interventional realisation, as C1 asks of a program's graph? Second, is there a contextual equivalence of concurrent programs, in the sense of \S\ref{sec:intv-wherewestand}, whose extensional content is a signalling structure, so that the pomset semantics of Kammar et al.\ and the signalling types of \S\ref{sec:hoqc} are the intensional and extensional faces of one notion?
\begin{itemize}[nosep]
\item[D1] \textit{(reading; position)} Pomsets and event structures from the interventional side: a position paper on what intervening on an event means and on the two questions above, informed by the threads theory of Kammar et al.\ \cite{kammar2026threads} and by Simmons's causal logic.
\end{itemize}

\subsubsection{Higher-order processes and quantum causality}\label{sec:prop-hoqc}

Higher-order quantum processes and my higher-order causal models are two generalisations of the same first-order object, in different directions. My dependent interventions live in a \emph{cartesian closed} category $[\mathbf{Bij},\mathcal C]$ generalising the base category $\mathcal C$; higher-order processes live in a \emph{compact closed} category generalising the \emph{Kleisli} category of the probability monad, the way $\CP$ generalises the deterministic processes $\mathrm{CPTP}$, and there is a classical version of every construction over $\Mat$ in which first-order processes are stochastic maps. The signalling notion of causality of Kissinger and Uijlen is defined on this compact closed side, and it is the interventional notion in disguise: a signalling constraint says that a local intervention at one node is undetectable at another. Three facts from the review line up with facts about my semantics. The signalling notion is strictly coarser than the causal structure, by fine tuning, exactly as $\equiv_2$ is coarser than dataflow. Quantum causal models are second-order processes, with nodes as slots, exactly as a program with intervention labels is a process with slots. And indefinite causal order is related to cyclic causal models \cite{barrett2021cyclic}, exactly as feedback is to fixpoints (\S\ref{sec:prop-expressivity}).

The precise form of the analogy is the following hypothesis. An intervention label is a \emph{slot} of a comb, and the intervener is the environment that fills it: $\dok(\ell\mapsto v)$ is the environment that discards what the program offers at the slot and returns $v$; no intervention is the environment that returns what it was offered; and, on $\Lk{2}$ extended with a first-order \emph{dependent} label kind $X\to P X$, the first-order shadow of the dependent intervention kind of $\LHO$, the contextual equivalence $\equiv_2^{\mathrm{dep}}$ under substitution of closed stochastic maps is equality of the combs obtained by leaving every label open, the slot being the $\mathsf{applyIntv}$ site. In the comb reading an environment is any stochastic map, and comb pairing is affine in it, which is why invariance transfers. Perfect interventions, the identity and the constant environments, affinely span the stochastic maps on a Boolean type (three points for a two-dimensional polytope) but not in general ($k+1$ points against a $k(k-1)$-dimensional polytope), so $\equiv_2$ with perfect interventions is strictly coarser than comb equality and perfect-intervention invariance strictly weaker than non-signalling from the slot. The POPL paper defines $\equiv_2$ for $\mathsf{Maybe}$-typed labels only, so $\equiv_2^{\mathrm{dep}}$ is a definition this item must supply. On the classical finite fragment this is a statement about stochastic maps and combs over $\Mat$, with no quantum theory in it, and it can be checked on paper with what I already know.
\begin{itemize}[nosep]
\item[D2] \textit{(medium risk)} The classical bridge: $\equiv_2^{\mathrm{dep}}$ defined, and a functor from finite $\Lk{2}$ denotations to combs over $\Mat$, in the $\CausC$ construction, with one slot per label, under which $\equiv_2^{\mathrm{dep}}$ becomes equality of combs and invariance under dependent interventions becomes non-signalling from the slot, with the perfect-intervention $\equiv_2$ characterised as the coarser equivalence it is. \emph{Done:} the functor, the two correspondence statements, and the separating example for perfect interventions. \emph{Why it may succeed:} it uses only Kissinger and Uijlen's first-order and one-way-signalling types over $\Mat$, which I know, and no quantum theory; it does not wait on the reading milestone M3.
\item[D3] \textit{(reading; stretch)} Quantum causal models \cite{barrett2019quantumcausalmodels,barrett2021cyclic} and Simmons's thesis in depth, towards the stretch goal: one category in which graph types, intervention grades and signalling types are three structures on the same objects, with the fine-tuning gap as the common theorem.
\end{itemize}

% ---------------------------------------------------------------------
\subsection{Evaluation, impact and risks}\label{sec:prop-risks}

\paragraph{Evaluation.} The thesis is theoretical with an executable component. Theorems are evaluated by the ``done'' criteria above: adequacy and, where a contextual equivalence is defined, full abstraction for each extended language (the POPL paper defines none for $\LHO$, so A1 is evaluated by differential testing of the library against the Haskell definitional implementation on the paper's examples), agreement theorems with the existing model-level notions (solution functions, causal spaces, abstraction, OmegaC's evaluation strategy), soundness of the algebraic theory of identification and admissibility of the do-calculus rules, and the correspondence statements of the classical bridge. The implementation is evaluated by reproducing the textbook causal-inference examples (mediation, instrumental variables, feedback) at realistic scale, by a comparison with ChiRho on an agreed benchmark set, and by the statement of which guarantees survive the host. The Haskell artifact \cite{ARTIFACT_wang_2026_23207221} is extended with every new construct so that every definition remains executable.

\paragraph{Impact.} If successful, the thesis gives denotational foundations for causal probabilistic programming, and hence for causal machine-learning languages, in which interventions are mathematical operations rather than evaluation strategies; the first program-level treatment of cyclic models, of abstraction and of identification; a practical library whose semantics is the specification of its implementation; and a precise account, on the classical fragment at least, of why signalling, interventional and concurrent causality are one notion with three intensional data.

\begin{table}[ht]
\centering\small
\begin{tabularx}{\textwidth}{llX}
\toprule
Item & Risk & Fallback\\
\midrule
A1 library & low & the examples without the benchmark comparison\\
A2 benchmark & medium & a qualitative comparison\\
A3 cycles, reading 1 & medium & $\mathbf{Set}$-models only; interventions on a fixed subset of the cycle variables\\
A4 cycles, reading 2 & high & a precise obstruction\\
B1 OmegaC, ChiRho & medium & OmegaC alone\\
B2 causal spaces & medium & the counterexample is itself a result\\
B3 abstraction & medium & the compiled-SCM case only\\
B4 synthetic & medium & representable Markov categories only\\
C1 graph types & medium & one of the two semantics\\
C2 algebraic theory & medium & soundness without the normalisation-box equations\\
C3 do-calculus & medium--high & admissibility of rules 1 and 2 only\\
D2 classical bridge & medium & one direction\\
D1, D3 & reading & a position chapter\\
\bottomrule
\end{tabularx}
\caption{Risks and fallbacks.}\label{tab:risks}
\end{table}

Table \ref{tab:risks} gives a fallback for each item. The generic risk is that the extensions are closer to existing lines than the base result was. The mitigation is the \emph{minimum viable thesis}: A1, A3, B1, B4 and C2, every item rated low or medium and each a program-level treatment of something that exists only at the model level; the originality lies in the program-level account, not in the choice of topics. The second risk is breadth, since the fourth direction reaches into two fields that \S\ref{sec:order} positions rather than surveys; the mitigation is that it is explicitly open-ended, that its deliverables are a position and one classical theorem, and that it can be cut to the classical bridge D2 without loss to the rest.

% ---------------------------------------------------------------------
\subsection{Plan of work, milestones, and training}\label{sec:plan}

I am in my fourth term. The baseline schedule is the minimum viable thesis; conditional items are marked. Reading milestones are M1 (the do-calculus, Markov conditions and d-separation; normalisation boxes; Bongers et al.\ and Forré--Mooij on cyclic models), M2 (Barrett, Lorenz and Oreshkov's quantum causal models, acyclic and cyclic), and M3 (Simmons's thesis, chapters 2--3; Kammar et al.\ on threads).

\begin{center}\small
\begin{tabularx}{\textwidth}{lX}
\toprule
Term & Milestones\\
\midrule
MT26 (4) & Transfer of status. POPL camera-ready. Reading M1.\\
HT27 (5) & A1 prototype. B1 (OmegaC). \emph{Deliverable:} note on the semantics of OmegaC and ChiRho.\\
TT27 (6) & A3. B2. C2 designed. Reading M2. \emph{Deliverable:} paper on cycles.\\
MT27--HT28 (7--8) & C2 soundness; C3 begun. B4. A2. \textbf{Confirmation of Status in HT28}, with the minimum viable thesis in hand.\\
TT28 (9) & Buffer for Confirmation of Status. C1. C3 completed. D2. \emph{Deliverable:} paper on the graph-centred account.\\
MT28 (10) & B3. Reading M3. A4 if A3 has landed, otherwise dropped.\\
HT29 (11) & D1, D3 as a position chapter. Consolidation.\\
TT29 (12) & Thesis writing only; submission by the end of TT29.\\
\bottomrule
\end{tabularx}
\end{center}
\noindent D2 depends on no reading milestone.

Dissemination targets are POPL, PLDI, LICS, MFPS and UAI. The thesis will be organised as: 1 Introduction; 2 Background (\S\ref{sec:foundations}--\ref{sec:synthesis} of this document); 3 Causal probabilistic programs: the hierarchy, contextual equivalence and the graded semantics \fillin{delineate your contribution from coauthors'}; 4 A practical language, and cycles; 5 The categorical constructions: relations and a synthetic account; 6 The graph-centred account and identification; 7 Causality in other fields: concurrency and higher-order processes; 8 Conclusion.

\paragraph{Training and engagement.} \fillin{required by the assessors' criteria: seminars attended (PL group, quantum group); courses audited; the PLanQC 2025 talk, with a paragraph on it; teaching; planned activities for the coming year}.

\subsection*{Statement of the subject of the thesis}

The semantics of causal probabilistic programming: contextual equivalence as the definition of interventional causal structure, made practical, extended to cycles and to a graph-centred account of identification, and related to signalling and concurrent causality.


\printbibliography

\end{document}
